% ENGR489 final report, working draft (IEEEtran).
% Text marked [TODO: ...] needs data or a fact only the author has.
% Text marked [CHECK: ...] is drafted from the repo but conflicts with, or is not
% confirmed by, another source. Resolve each before submission.
\documentclass[conference]{IEEEtran}
\usepackage{amsmath,graphicx,booktabs,siunitx,cite,url}
\sisetup{detect-all}

\begin{document}

\title{Design of a Robotic Platform for Magnetic Field Mapping in a Helmholtz Coil System}
\author{\IEEEauthorblockN{Noah-S\={e}a Nui}
\IEEEauthorblockA{Supervised by Christopher Hollitt\\
School of Engineering and Computer Science, Victoria University of Wellington}}
\maketitle

\begin{abstract}
Mapping the field of a Helmholtz cage by hand is slow and its results depend on the
operator. This project built a delta robot to position a three-axis magnetometer
automatically. The arms are 3D-printed PLA and carbon-fibre tube, and the probe sits
on a PLA holder, so no ferromagnetic material is near the sensor. Three hobby servos
move the probe to the points of a Cartesian grid through a closed-form inverse
kinematic solution. A host program moves, waits, averages 20 readings and logs each
point. With a position correction fitted to the as-built robot, the corner error of a
$100\times 100$\,mm scan box fell from about \SI{19}{mm} to \SI{2.3}{mm} rms, and
repeatability was \SIrange{2}{2.5}{mm}. This does not meet the \SI{1}{mm}
requirement. A 75-point scan takes about \SI{14}{min}. Repeated scans of a fixed
magnet agreed to \SI{0.32}{\micro T} rms in $|B|$.
[TODO: finish with the in-cage robot-interference and field-vs-theory results once measured.]
\end{abstract}

\begin{IEEEkeywords}
Delta robot, magnetic field mapping, Helmholtz coil, parallel kinematics, magnetometer.
\end{IEEEkeywords}

\section{Introduction}
\subsection{Motivation}
Helmholtz cages produce a nearly uniform field and are used to calibrate magnetic
sensors and to test CubeSat attitude determination and control systems.
The VUW three-axis cage was built in 2025 [TODO: cite the 2025 report].
Its uniformity was measured by hand on a \SI{2}{cm} grid over three planes, which took
about one hour, and probe misorientation was identified as a source of error.
That report recommended an automated, motorised mapping system as future work.
This project delivers that system.

\subsection{Problem Statement}
Manual mapping limits resolution, repeatability and orientation consistency. A
mapping system must also not disturb the field it measures, so every part near the
probe has to be non-ferromagnetic.

\subsection{Requirements}
The platform was required to:
\begin{itemize}
\item position a three-axis magnetometer with accuracy of the order of \SI{1}{mm}
within a $100\,\mathrm{mm}\times100\,\mathrm{mm}\times300\,\mathrm{mm}$ test volume;
\item measure the field vector $(B_x,B_y,B_z)$ at discrete grid points;
\item provide an adjustable scan resolution;
\item minimise magnetic interference introduced by the platform; and
\item produce repeatable measurements across multiple scans.
\end{itemize}
The project budget was \$500. [TODO: state the safety requirement if one was set; the
firmware has no hardware E-stop, only the servo-power switch and a software disable.]

\subsection{Summary of Outcomes}
A delta robot was built and driven by an Arduino UNO R4 Minima. Grid scanning,
adjustable resolution, non-magnetic construction and scan repeatability were achieved.
Positioning accuracy and the \SI{300}{mm} depth were only partly met: after correction
the robot reached the full $100\times299$\,mm box at $z=\SI{-650}{mm}$ with a corner
error of \SI{2.3}{mm} rms against the \SI{1}{mm} target. [CHECK: the draft earlier quoted
a reachable depth of \SI{276}{mm}; the firmware README gives $z\approx-550$ to
\SI{-793}{mm} on axis and the correction validation gives \SI{299}{mm}. Use one
measured figure.] [TODO: headline in-cage numbers.]

\subsection{Sustainability}
The work relates to UN SDG 9 (industry, innovation and infrastructure) because it
gives a low-cost automated alternative to commercial mapping equipment
[TODO: cite UN SDGs]. It relates to SDG 4 (quality education) if the platform is reused in teaching labs. The other goals are
not materially affected by a laboratory instrument of this size.

\begin{table}[t]
\caption{Performance metrics and their relevance to the project requirements.}
\label{tab:metrics}
\centering\small
\begin{tabular}{p{0.38\columnwidth}p{0.52\columnwidth}}
\toprule
Metric & Justification \\
\midrule
Positioning accuracy (mm) & Direct measure of the \SI{1}{mm} requirement \\
Positioning repeatability (mm) & Separates random error from correctable bias \\
Reachable workspace (mm) & Must cover the $100\times100\times300$\,mm volume \\
Robot-induced field disturbance (nT) & Measures the non-interference requirement \\
Field error vs theory, RMS and max (nT) & Checks the map against Eq.~(1) \\
Scan repeatability (nT) & Measures the repeatability requirement \\
Scan time per point (s) & Compared against the manual baseline \\
Total cost (NZD) & Must stay within \$500 \\
\bottomrule
\end{tabular}
\end{table}

\section{Background and Related Work}
\subsection{Helmholtz Coil Fields}
[TODO: Present the square Helmholtz axial-field equation with its source, define each
variable with units, and tabulate the VUW cage parameters (side length, separation,
turns). Confirm the cage is unchanged since 2025.]
\begin{equation}
B_z = [\mathrm{TODO:\ equation}]
\end{equation}

\subsection{Existing Field Mapping Approaches}
[TODO: Manual probing; Cartesian gantry mappers; serial robot arms; commercial
Hall-probe mapping systems. For each give advantages, disadvantages, magnetic
compatibility and cost. Cite real sources only.]

\subsection{Delta Robots}
The delta robot was invented by Clavel [1]. Three rotary actuators on a fixed base
drive three parallel arm pairs that hold a small platform, which keeps the
platform level and gives it three translational degrees of freedom. Deltas are
common in pick-and-place machines and some 3D printers because the moving mass is
low. The reachable workspace is a rounded volume, not a box.

\subsection{Novelty of This Work}
The platform is a delta robot built from PLA and carbon-fibre tube with no metal near
the probe, designed for the VUW cage and costing a fraction of a commercial mapper.
[TODO: confirm the cost figure and state clearly what is new relative to the sources
in Section II-B.]

\subsection{Performance Metrics}
Table~\ref{tab:metrics} lists the metrics used in Section~VI.

\section{Tools and Methodology}
SolidWorks was used for the assembly, with mates for the ball joints, DXF export for
laser cutting and STL export for printing. Printed holes were enlarged by
\SI{0.2}{mm} so that fasteners and rod ends fit. Kinematics were simulated in MATLAB
and Python, adapted from the open-source project of Wrona and Radziwi\l{}\l{}ko [2].
The firmware targets an Arduino UNO R4 Minima and was built with the Arduino IDE.
The magnetometer is a Phidget 1044 (PhidgetSpatial Precision 3/3/3, serial 302277),
read through the Phidget22 API. Code is under Git.
[TODO: GitLab link.]

\begin{figure}[t]
\centering
[TODO: architecture diagram]
\caption{System architecture of the field mapping platform.}
\end{figure}

\section{Design}
\subsection{System Architecture}
The host PC runs the scan program. It sends framed JSON commands over USB CDC serial
to the R4 Minima, which generates PWM on D9, D10 and D11 for the three servos. The
firmware streams a status line every \SI{20}{ms}. The servos run from a separate 6\,V
AA pack with a common ground. The Phidget 1044 connects to the host directly over USB,
not through the Arduino. [TODO: potentiometer feedback on A0; it is not used by the v1
firmware, which is open loop.]

\subsection{Revised Architecture Selection}
The project proposal selected a three-axis Cartesian gantry, on the basis that its
decoupled linear axes give a direct relationship between actuator displacement and
probe position. Parallel robots, including the delta configuration, were rejected at
that stage because of their coupled, nonlinear kinematics, their non-rectangular
workspace, and the presence of linkages around the measurement volume.

During detailed design, the gantry could not be delivered within the ENGR489 budget of
\$500. The aluminium support structure needed to hold three axes rigidly above the
Helmholtz coil, together with the linear rails, belts, pulleys and lead screw, brought
the costed design to approximately \$[TODO: X]. Several redesigns were attempted,
including [TODO: e.g.\ shorter travel, lower-cost rails], but none brought the design
within budget without compromising rigidity.

A delta robot was therefore adopted. It uses three rotary actuators fixed to a
stationary base, so the three Futaba S3306MG servos already available to the project
could drive the arms directly, with no rotary-to-linear conversion. The arms were
built from 3D-printed PLA and carbon-fibre tube. The final cost of the platform was
\$[TODO: Y]. The proposal's objections were addressed as follows:
\begin{itemize}
\item \emph{Nonlinear kinematics:} an inverse kinematic solution adapted from [2]
converts Cartesian targets into servo angles (Section IV-C), so scans are still
defined on a Cartesian grid.
\item \emph{Non-rectangular workspace:} the reachable workspace was analysed for the
as-built geometry (Section IV-D). [CHECK: depth figure, see Section I-D.] The shortfall
against the \SI{300}{mm} target is a direct consequence of the change in architecture
and is discussed in Section VI-G.
\item \emph{Linkages near the measurement volume:} the effect of the robot on the field
was to be quantified experimentally (Section VI-D).
\end{itemize}
The change traded a simpler kinematic model and a guaranteed rectangular workspace for
a design that was affordable and reused existing actuators.

\begin{table}[t]
\caption{As-built geometric parameters of the delta robot.}
\label{tab:geom}
\centering\small
\begin{tabular}{lr}
\toprule
Parameter & Value (mm) \\
\midrule
Base triangle side, $s_b$ & 175 \\
Platform triangle side, $s_p$ & 75 \\
Upper arm length, $L_1$ & 177 \\
Forearm length, $L_2$ & 625 \\
Probe offset below platform joint centre & 21 \\
\bottomrule
\end{tabular}
\end{table}
[CHECK: values are from the calibrated firmware recorded in the scan metadata
(branch \texttt{claude/amazing-lovelace-onixrv}). The \texttt{engr489-firmware}
copy on this branch still has $s_p=150$, $L_1=180$. Project notes list $s_b=450$,
$s_p=18$, $L_1=180$, $L_2=600$. Confirm against the physical robot and keep one table.]

\subsection{Kinematics}
Each arm $i$ is a bicep of length $L_1$ rotating about a horizontal axis on the base
triangle, and a forearm of length $L_2$ joining the elbow to the platform. For a platform
position $(x,y,z)$ the closed loop gives, per arm,
\begin{equation}
E_i\cos\theta_i + F_i\sin\theta_i + G_i = 0,
\end{equation}
with $E_i$, $F_i$, $G_i$ functions of $(x,y,z)$ and the geometry of Table~\ref{tab:geom}.
Substituting $t=\tan(\theta_i/2)$ gives
\begin{equation}
\theta_i = 2\arctan\frac{-F_i\pm\sqrt{E_i^2+F_i^2-G_i^2}}{G_i-E_i}.
\end{equation}
The firmware tries the $+$ root first and the $-$ root if the first is outside the
servo limits, and rejects the point if neither is real and in range. [TODO: write
out $E_i,F_i,G_i$ from \texttt{kinematics.cpp} with units, and define $\theta$ sign
convention: positive is bicep down.] Servo pulse width is linear in angle,
$\mu_i = c_i + k_i\theta_i$, with $c_i$ the pulse for a level bicep and $k_i$ in
\SI{}{\micro s/deg}.

\subsection{Workspace Analysis}
Reach was found by evaluating the inverse kinematics over a grid. Limits come from
the servo angle range of $-20^\circ$ to $70^\circ$ (set by horn fouling and the plate
slot, measured with all three arms enabled; single-arm limits were falsely tight because
the plate sagged), the real-root condition and the frame window.
[TODO: plot of reachable volume against the target box, and the $s_b/L$ trade-offs.]

\subsection{Mechanical Design}
[TODO: Top plate (acrylic prototype, aluminium final); servo mounting and arm slots;
bicep box section; elbow yoke in double shear; carbon-fibre forearms bonded into bored
rod ends; bonding jig; probe holder; stand of 2020 extrusion legs at $120^\circ$.]

\subsection{Material Selection and Magnetic Considerations}
The non-ferromagnetic rule was the main design driver. [TODO: table of materials by
component.] Moving parts near the probe are PLA, carbon fibre, nylon and aluminium.
The known compromise is the SIL6T/K M6 rod ends, whose balls are steel. They were
accepted for this build and should be replaced by plastic or non-ferrous joints.
The servos contain magnets; their effect appears in the baseline scan (Section VI-D).

\subsection{Electronics}
The UNO R4 Minima was chosen over the R4 WiFi to avoid RF emission near the
magnetometer. Servo power is separate from logic power, with a shared ground.
Brown-out of the 4$\times$AA supply under servo load is a known failure mode.
[TODO: potentiometer feedback modification and voltage range; E-stop.]

\subsection{Sustainability Considerations}
[TODO: Environmental: PLA and the disposable AA pack (recommend rechargeable).
Economic: cost against the gantry and commercial systems. Technical: off-the-shelf
parts, open-source code, replaceable printed parts. Social: argue relevance or not.]

\section{Implementation}
\subsection{Mechanical Build}
[TODO: Fabrication steps, photos, problems met and fixes.]

\subsection{Firmware}
The sketch is split into host-testable modules: kinematics, a frame parser, and a
motion planner, with only the \texttt{.ino} file touching \texttt{Servo} and
\texttt{Serial}. It keeps the serial protocol of [2] byte-for-byte so the upstream
Python GUI still works, and replaces the stepper control with servo pulses clamped to
\SIrange{830}{2170}{\micro s}. [CHECK: the draft said 800--2200\,$\mu$s and EEPROM
storage; \texttt{config.h} has 830--2170\,$\mu$s and the calibration is compiled in,
EEPROM is listed as future work.] Moves are planned with a smoothstep profile at a tick
of \SI{20}{ms}, and a status field reports errors such as an unreachable target. Ten
host-side test groups cover kinematics, parity between the C++ and Python solutions,
the parser and the planner. [TODO: GitLab link.]

\subsection{Servo Calibration}
Each servo was calibrated with an angle finder on the bicep. The pulse width at
level was found first, then a five-point linear fit at 1200, 1360, 1520, 1680 and
\SI{1840}{\micro s} gave the slope. Angle limits were then found with all three arms
enabled. Table~\ref{tab:cal} gives the final values.
\begin{table}[t]
\caption{Servo calibration (calibrated firmware).}
\label{tab:cal}
\centering\small
\begin{tabular}{lccc}
\toprule
Arm (pin) & $c_i$ ($\mu$s) & $k_i$ ($\mu$s/deg) & Limits (deg) \\
\midrule
1 (D9)  & 1456 & 9.8315 & $-20$ to 70 \\
2 (D10) & 1373 & 9.7664 & $-20$ to 70 \\
3 (D11) & 1393 & 9.6658 & $-20$ to 70 \\
\bottomrule
\end{tabular}
\end{table}
[TODO: fit residuals and the exit criteria used.]

\subsection{Host Software and Scan Routine}
\texttt{field\_scan.py} opens the magnetometer, then the robot, parks at
$(0,0,-650)$\,mm, and visits the grid in serpentine order. At each point it sends a
move, waits for the robot to report that it has stopped, settles for \SI{5}{s},
averages 20 readings taken \SI{20}{ms} apart, and appends one CSV row. Unreachable
points are logged as NaN. Each scan writes a metadata file with every setting,
the sensor range and the firmware geometry. MATLAB scripts subtract a
baseline scan matched by position and plot layers, side views and 3D arrows.
The grid is defined in probe coordinates; the commanded platform position is the
probe position less the \SI{21}{mm} offset. A \SI{0.5}{s} settle was tried first and
replaced by \SI{5}{s} because the probe was still moving.
[TODO: confirm the reason for the change and describe the CSV columns in a table.]

\section{Evaluation}
\subsection{Test Setup}
All scans so far were run with the robot outside the cage, coils off, a strong
permanent magnet near the probe as the field source, and a no-magnet baseline.
Each scan covers $5\times5\times3$ points at $x,y\in[-50,50]$\,mm and
$z\in[-675,-625]$\,mm (75 points). [TODO: photo; in-cage test conditions.]

\subsection{Positioning Accuracy and Repeatability}
Position was checked with a pen and ruler over $x\pm50$, $y\pm150$ and
$z=-600$ to \SI{-700}{mm}. Without correction the servos give way under load off
centre and the corners of the box land about \SI{19}{mm} inside the target. The
\texttt{hybrid} correction predicts each servo's shortfall from the Jacobian and aims
past the target by that amount. It reduced the corner error to \SI{2.3}{mm} rms,
with about \SI{5}{mm} left at the box centre at $z=-650$\,mm caused by platform tilt.
Repeatability was \SIrange{2}{2.5}{mm}. The \SI{1}{mm} accuracy requirement was
therefore not met. [TODO: tabulate commanded vs measured points and a return-to-point
test; cite the calibration notes.]

\subsection{Achieved Workspace}
With the corrected limits the on-axis reach is $z\approx-550$ to \SI{-793}{mm}, and
about 89\,\% of a $\pm150$\,mm by $z\in[-790,-550]$\,mm grid is reachable. The
corrected scan box is $100\times299$\,mm at $z=-650$. [TODO: measured against
predicted reach; resolve the \SI{276}{mm} figure.]

\subsection{Robot-Induced Field Disturbance}
Not yet measured in the cage. [TODO: robot present vs absent, powered vs unpowered,
against the sensor noise floor.] In the bench scans, the mean reading with no magnet
was \SI{72.0}{\micro T}, close to the expected Earth field of \SIrange{55}{60}{\micro T}
plus the contribution of the servo magnets and sensor offset [TODO: attribute the
\SI{12}{\micro T} excess]. The per-point noise of the 20-reading average was
\SIrange{0.06}{0.10}{\micro T} per axis, which sets the noise floor.

\subsection{Field Map and Comparison with Theory}
[TODO: Field maps; RMS and max error vs Eq.~(1); uniformity as percent deviation
from the centre value. Requires coil-on scans in the cage.]

\subsection{Scan Repeatability and Scan Time}
Table~\ref{tab:rep} gives the difference in $|B|$ between repeated scans of the same
grid with a fixed magnet. The uncorrected pair differed by \SI{0.32}{\micro T} rms
(\SI{0.89}{\micro T} max) and the corrected pair by \SI{0.28}{\micro T} rms
(\SI{0.74}{\micro T} max). The field varies from \SI{73}{\micro T} to
\SI{87}{\micro T} over the grid, so the repeat error is about 2\,\% of the span.
A scan of 75 points took \SIrange{830}{862}{s}, or \SI{11.1}{s} to \SI{11.5}{s} per
point, dominated by the \SI{5}{s} settle. [TODO: compare with the one-hour manual
baseline in the 2025 report, and give the number of points in that scan.]
\begin{table}[t]
\caption{Difference in $|B|$ between repeated 75-point scans (strong magnet, coils off).}
\label{tab:rep}
\centering\small
\begin{tabular}{lcc}
\toprule
Pair & RMS ($\mu$T) & Max ($\mu$T) \\
\midrule
No correction, run 1 vs 2 & 0.32 & 0.89 \\
Hybrid correction, run 1 vs 2 & 0.28 & 0.74 \\
\bottomrule
\end{tabular}
\end{table}

\subsection{Limitations and Future Improvements}
\begin{itemize}
\item \emph{Accuracy:} \SI{2.3}{mm} rms against \SI{1}{mm}. The servos deflect under
load and the platform tilts; stiffer servos or closed-loop control on the
potentiometers would help.
\item \emph{Workspace:} geometry change to reach the full \SI{300}{mm} depth. [TODO]
\item \emph{Rod ends:} replace the steel-ball rod ends.
\item \emph{Servo resolution and deadband,} and open-loop control.
\item \emph{Supply:} AA pack sags under load.
\item \emph{Sensor noise} and the untested in-cage disturbance.
\end{itemize}

\section{Discussion}
[TODO: Were the design choices right in hindsight? Was the proposal's view of delta
robots correct? Sustainability impact. Note that the proposal's concern about coupled
kinematics was handled in software, but the concern about positioning accuracy
appeared in a different form: servo compliance, which a gantry would not have.]

\section{Conclusion}
[TODO: Objectives vs outcomes, one short paragraph each.]

\begin{thebibliography}{1}
\bibitem{clavel} R. Clavel, ``Device for the movement and positioning of an element in space,'' U.S.\ Patent 4\,976\,582, Dec.\ 11, 1990.
\bibitem{delta} J. Wrona and G. Radziwi\l{}\l{}ko, ``Delta-Robot,'' GitHub repository, commit 48a8038. [Online]. Available: https://github.com/grzesiek2201/Delta-Robot
\end{thebibliography}

\end{document}
